\documentclass[uplatex,a4paper,11pt]{jsarticle}
\usepackage[margin=22mm]{geometry}
\usepackage[dvipdfmx]{hyperref,xcolor}
\usepackage{fvextra,enumitem,longtable,booktabs,amssymb}
\hypersetup{colorlinks=true,linkcolor=blue!55!black,urlcolor=blue!65!black,
pdftitle={Raspberry Pi TeXNote版組サーバー構築手順}}
\definecolor{commandbg}{RGB}{244,246,248}
\DefineVerbatimEnvironment{command}{Verbatim}{fontsize=\small,frame=single,
rulecolor=\color{gray!55},breaklines=true,breakanywhere=true}
\setlist{nosep}
\title{Raspberry Pi TeXNote版組サーバー構築手順\\
\large 状態確認からpdfLaTeX/CJK用Type 1フォント導入まで}
\author{}\date{2026年8月15日}
\usepackage{comment}
\usepackage{listings, jlisting}
\lstset{
	%プログラム言語(複数の言語に対応，C,C++も可)
	language = Python,
	%背景色と透過度
	backgroundcolor={\color[gray]{.90}},
	%枠外に行った時の自動改行
	breaklines = true,
	%自動改行後のインデント量(デフォルトでは20[pt])
	breakindent = 10pt,
	%標準の書体
	basicstyle = \ttfamily\scriptsize,
	%コメントの書体
	commentstyle = {\itshape \color[cmyk]{1,0.4,1,0}},
	%関数名等の色の設定
	classoffset = 0,
	%キーワード(int, ifなど)の書体
	keywordstyle = {\bfseries \color[cmyk]{0,1,0,0}},
	%表示する文字の書体
	stringstyle = {\ttfamily \color[rgb]{0,0,1}},
	%枠 "t"は上に線を記載, "T"は上に二重線を記載
	%他オプション：leftline，topline，bottomline，lines，single，shadowbox
	frame = TBrl,
	%frameまでの間隔(行番号とプログラムの間)
	framesep = 5pt,
	%行番号の位置
	numbers = left,
	%行番号の間隔
	stepnumber = 1,
	%行番号の書体
	numberstyle = \tiny,
	%タブの大きさ
	tabsize = 4,
	%キャプションの場所("tb"ならば上下両方に記載)
	captionpos = t
}

\begin{document}
\maketitle
\begin{abstract}
iPad版・iPhone版TeXNoteから、同一LAN上のRaspberry PiへTeX文書を送り、
PDFを生成するサーバーを再構築するための手順書である。
Debian 12、TeX Live、FastAPI/Uvicorn、APIトークン、systemdによる自動起動、
およびpdfLaTeX/CJKで必要になるType 1フォント一式までを扱う。
\end{abstract}
\tableofcontents
\newpage

\section{完成時の構成}

実際に構築した環境は次のとおりである。再構築時に異なる値は読み替える。
\begin{longtable}{@{}ll@{}}
\toprule 項目 & 値\\ \midrule
機種 & Raspberry Pi 5 Model B Rev 1.0\\
OS & Debian GNU/Linux 12 (bookworm)\\
CPU・メモリ & aarch64・約8\,GiB\\
ストレージ & 約29\,GB\\
ホスト名・ユーザー & raspberrypi・mat\\
LAN内IPアドレス & 192.168.3.11\\
サーバーディレクトリ & /home/mat/texnote-server\\
待受ポート & TCP 8000\\
TeXエンジン & LuaLaTeX、XeLaTeX、pdfLaTeX、upLaTeX、pLaTeX\\
\bottomrule
\end{longtable}
コマンドは「Raspberry Pi」または「Mac」と明記した側で実行する。

%command1.sh

\section{Raspberry Piの状態確認}
Macから接続する。
\begin{command}
ssh mat@192.168.3.11
\end{command}
Raspberry Pi上で状態を確認する。
\begin{command}
cat /etc/os-release
uname -m
df -h hostname -I
tr -d '\0' </proc/device-tree/model
echo
free -h
hostnamectl
timedatectl
sudo apt update
apt list --upgradable
test -f /var/run/reboot-required \
  && cat /var/run/reboot-required \
  || echo "再起動は要求されていません"
\end{command}

\begin{comment}
command1.sh の結果例

\begin{verbatim}
PRETTY_NAME="Debian GNU/Linux 12 (bookworm)"
NAME="Debian GNU/Linux"
VERSION_ID="12"
VERSION="12 (bookworm)"
VERSION_CODENAME=bookworm
ID=debian
HOME_URL="https://www.debian.org/"
SUPPORT_URL="https://www.debian.org/support"
BUG_REPORT_URL="https://bugs.debian.org/"
aarch64
df: invalid option -- 'I'
Try 'df --help' for more information.
Raspberry Pi 5 Model B Rev 1.0
               total        used        free      shared  buff/cache   available
Mem:           7.9Gi       563Mi       6.7Gi        58Mi       731Mi       7.3Gi
Swap:          199Mi          0B       199Mi
 Static hostname: raspberrypi
       Icon name: computer
      Machine ID: 34fa9c445800446dadcbd0259e955b64
         Boot ID: f1423dad035949168771f530b7bac316
Operating System: Debian GNU/Linux 12 (bookworm)
          Kernel: Linux 6.12.96+rpt-rpi-2712
    Architecture: arm64
               Local time: Fri 2026-08-14 22:43:00 JST
           Universal time: Fri 2026-08-14 13:43:00 UTC
                 RTC time: Fri 2026-08-14 13:43:00
                Time zone: Asia/Tokyo (JST, +0900)
System clock synchronized: yes
              NTP service: active
          RTC in local TZ: no
Hit:1 http://deb.debian.org/debian bookworm InRelease
Hit:2 http://deb.debian.org/debian-security bookworm-security InRelease
Hit:3 http://deb.debian.org/debian bookworm-updates InRelease
Hit:4 http://archive.raspberrypi.com/debian bookworm InRelease
Reading package lists... Done
Building dependency tree... Done
Reading state information... Done
All packages are up to date.
Listing... Done
再起動は要求されていません
\end{verbatim}
\end{comment}

画像ファイルがepsであった場合、gs（ghostscript）が必要になる
\begin{command}
sudo apt install ghostscript
command -v gs
gs --version
\end{command}

概ね次の様に出れば問題なくインストールできている
\begin{verbatim}
/usr/bin/gs
10.00.0
\end{verbatim}

必要なら更新して再起動する。
\begin{command}
sudo apt full-upgrade
sudo reboot
\end{command}
IPが変わるとアプリから接続できなくなるため、ルーターのDHCP予約を推奨する。

\section{TeX Liveと日本語環境}

まず導入をシミュレーションする。
\begin{command}
sudo apt-get --simulate install \
  texlive \
  texlive-latex-extra \
  texlive-luatex \
  texlive-xetex \
  texlive-lang-japanese \
  texlive-pictures \
  texlive-science \
  fonts-noto-cjk
\end{command}
大量の削除がないことを確認後、実際に導入する。
\begin{command}
sudo apt install \
  texlive \
  texlive-latex-extra \
  texlive-luatex \
  texlive-xetex \
  texlive-lang-japanese \
  texlive-pictures \
  texlive-science \
  fonts-noto-cjk
\end{command}
エンジンを確認する。

%command2.sh

\begin{command}
command -v lualatex
command -v xelatex
command -v pdflatex
command -v uplatex
command -v platex
command -v dvipdfmx
command -v extractbb
lualatex --version | head -n 1
xelatex --version | head -n 1
pdflatex --version | head -n 1
uplatex --version | head -n 1
platex --version | head -n 1
dvipdfmx --version | head -n 1
\end{command}

\begin{comment}
command2.sh 実行結果例

\begin{verbatim}
/usr/bin/lualatex
/usr/bin/xelatex
/usr/bin/pdflatex
/usr/bin/uplatex
/usr/bin/platex
/usr/bin/dvipdfmx
/usr/bin/extractbb
This is LuaHBTeX, Version 1.15.0 (TeX Live 2022/Debian)
XeTeX 3.141592653-2.6-0.999994 (TeX Live 2022/Debian)
pdfTeX 3.141592653-2.6-1.40.24 (TeX Live 2022/Debian)
e-upTeX 3.141592653-p4.0.0-u1.28-220214-2.6 (utf8.uptex) (TeX Live 2022/Debian)
e-pTeX 3.141592653-p4.0.0-220214-2.6 (utf8.euc) (TeX Live 2022/Debian)
This is dvipdfmx Version 20211117 by the DVIPDFMx project team,
\end{verbatim}
\end{comment}

\section{5種類のエンジンを単体テスト}

%command3.sh

\begin{command}
mkdir -p ~/tex-tests
cd ~/tex-tests

cat > lualatex-test.tex <<'EOF'
\documentclass{ltjsarticle}
\begin{document}
LuaLaTeX日本語テスト。\quad $e^{i\pi}+1=0$
\end{document}
EOF
lualatex -interaction=nonstopmode -halt-on-error lualatex-test.tex

cat > xelatex-test.tex <<'EOF'
\documentclass[xelatex,ja=standard]{bxjsarticle}
\begin{document}
XeLaTeX日本語テスト。\quad $e^{i\pi}+1=0$
\end{document}
EOF
xelatex -interaction=nonstopmode -halt-on-error xelatex-test.tex

cat > pdflatex-test.tex <<'EOF'
\documentclass{article}
\begin{document}
pdfLaTeX test.\quad $e^{i\pi}+1=0$
\end{document}
EOF
pdflatex -interaction=nonstopmode -halt-on-error pdflatex-test.tex

cat > uplatex-test.tex <<'EOF'
\documentclass[uplatex,dvipdfmx]{jsarticle}
\begin{document}
upLaTeX日本語テスト。\quad $e^{i\pi}+1=0$
\end{document}
EOF
uplatex -interaction=nonstopmode -halt-on-error uplatex-test.tex
dvipdfmx uplatex-test.dvi

cat > platex-test.tex <<'EOF'
\documentclass[dvipdfmx]{jsarticle}
\begin{document}
pLaTeX日本語テスト。\quad $e^{i\pi}+1=0$
\end{document}
EOF
platex -interaction=nonstopmode -halt-on-error platex-test.tex
dvipdfmx platex-test.dvi

ls -lh *-test.pdf
file *-test.pdf
\end{command}
pdfLaTeXはまず欧文と数式で確認し、CJK日本語用フォントは後節で導入する。

\begin{comment}
command3.sh 実行結果

\begin{verbatim}
This is LuaHBTeX, Version 1.15.0 (TeX Live 2022/Debian)
 restricted system commands enabled.
(./lualatex-test.tex
LaTeX2e <2022-11-01> patch level 1
 L3 programming layer <2023-01-16>
(/usr/share/texlive/texmf-dist/tex/luatex/luatexja/ltjsarticle.cls
Document Class: ltjsarticle 2022/09/12 ltjsclasses
(/usr/share/texlive/texmf-dist/tex/luatex/luatexja/luatexja.sty
(/usr/share/texlive/texmf-dist/tex/luatex/luatexja/luatexja-core.sty
(/usr/share/texlive/texmf-dist/tex/luatex/luatexbase/luatexbase.sty
(/usr/share/texlive/texmf-dist/tex/luatex/ctablestack/ctablestack.sty))
(/usr/share/texlive/texmf-dist/tex/generic/ltxcmds/ltxcmds.sty)
(/usr/share/texlive/texmf-dist/tex/generic/pdftexcmds/pdftexcmds.sty
(/usr/share/texlive/texmf-dist/tex/generic/infwarerr/infwarerr.sty)
(/usr/share/texlive/texmf-dist/tex/generic/iftex/iftex.sty))
(/usr/share/texlive/texmf-dist/tex/latex/xkeyval/xkeyval.sty
(/usr/share/texlive/texmf-dist/tex/generic/xkeyval/xkeyval.tex
(/usr/share/texlive/texmf-dist/tex/generic/xkeyval/xkvutils.tex
(/usr/share/texlive/texmf-dist/tex/generic/xkeyval/keyval.tex))))
(/usr/share/texlive/texmf-dist/tex/latex/etoolbox/etoolbox.sty)
(/usr/share/texlive/texmf-dist/tex/latex/everyhook/everyhook.sty
(/usr/share/texlive/texmf-dist/tex/latex/svn-prov/svn-prov.sty))
(/usr/share/texlive/texmf-dist/tex/luatex/luatexja/ltj-base.sty)
(/usr/share/texlive/texmf-dist/tex/luatex/luatexja/ltj-kinsoku.tex)
(/usr/share/texlive/texmf-dist/tex/luatex/luatexja/ltj-latex.sty
(/usr/share/texlive/texmf-dist/tex/luatex/luatexja/patches/lltjfont.sty
(/usr/share/texlive/texmf-dist/tex/latex/base/tuenc.def))
(/usr/share/texlive/texmf-dist/tex/luatex/luatexja/patches/lltjdefs.sty
luaotfload | db : Font names database not found, generating new one.
luaotfload | db : This can take several minutes; please be patient.
(/usr/share/texlive/texmf-dist/tex/luatex/luatexja/jfm-ujisv.lua)
(/usr/share/texlive/texmf-dist/tex/luatex/luatexja/jfm-ujis.lua))
(/usr/share/texlive/texmf-dist/tex/luatex/luatexja/patches/lltjcore.sty
(/usr/share/texlive/texmf-dist/tex/latex/l3kernel/expl3.sty
(/usr/share/texlive/texmf-dist/tex/latex/l3backend/l3backend-luatex.def)))
(/usr/share/texlive/texmf-dist/tex/luatex/luatexja/patches/lltjp-atbegshi.sty)
(/usr/share/texlive/texmf-dist/tex/luatex/luatexja/patches/lltjp-geometry.sty
(/usr/share/texlive/texmf-dist/tex/generic/iftex/ifluatex.sty))))
(/usr/share/texlive/texmf-dist/tex/luatex/luatexja/luatexja-compat.sty))
(/usr/share/texlive/texmf-dist/tex/platex/jsclasses/jslogo.sty)
(/usr/share/texlive/texmf-dist/tex/latex/sttools/stfloats.sty)
(/usr/share/texlive/texmf-dist/tex/luatex/luatexja/patches/lltjp-stfloats.sty))

No file lualatex-test.aux.
(/usr/share/texlive/texmf-dist/tex/latex/base/ts1cmr.fd) [1{/var/lib/texmf/font
s/map/pdftex/updmap/pdftex.map}] (./lualatex-test.aux))
 792 words of node memory still in use:
   7 hlist, 1 vlist, 4 rule, 2 glue, 3 kern, 1 glyph, 129 attribute, 49 glue_sp
ec, 20 attribute_list, 1 write, 9 user_defined nodes
   avail lists: 1:7,2:952,3:12,4:17,5:24,6:3,7:106,8:7,9:34,11:4
</usr/share/texlive/texmf-dist/fonts/opentype/public/haranoaji/HaranoAjiMincho-
Regular.otf></usr/share/texmf/fonts/opentype/public/lm/lmroman10-regular.otf></
usr/share/texlive/texmf-dist/fonts/type1/public/amsfonts/cm/cmmi10.pfb></usr/sh
are/texlive/texmf-dist/fonts/type1/public/amsfonts/cm/cmmi7.pfb></usr/share/tex
live/texmf-dist/fonts/type1/public/amsfonts/cm/cmr10.pfb>
Output written on lualatex-test.pdf (1 page, 28651 bytes).
Transcript written on lualatex-test.log.
This is XeTeX, Version 3.141592653-2.6-0.999994 (TeX Live 2022/Debian) (preloaded format=xelatex)
 restricted \write18 enabled.
entering extended mode
(./xelatex-test.tex
LaTeX2e <2022-11-01> patch level 1
L3 programming layer <2023-01-16>
(/usr/share/texlive/texmf-dist/tex/latex/bxjscls/bxjsarticle.cls
Document Class: bxjsarticle 2022/04/10 v2.7a BXJS document classes
(/usr/share/texlive/texmf-dist/tex/latex/tools/calc.sty)
(/usr/share/texlive/texmf-dist/tex/latex/graphics/keyval.sty)
(/usr/share/texlive/texmf-dist/tex/latex/bxjscls/bxjscompat.sty)
(/usr/share/texlive/texmf-dist/tex/generic/iftex/ifpdf.sty
(/usr/share/texlive/texmf-dist/tex/generic/iftex/iftex.sty))
(/usr/share/texlive/texmf-dist/tex/latex/geometry/geometry.sty
(/usr/share/texlive/texmf-dist/tex/generic/iftex/ifvtex.sty))
(/usr/share/texlive/texmf-dist/tex/latex/bxwareki/bxwareki.sty
(/usr/share/texlive/texmf-dist/tex/latex/bxwareki/bxwareki-cd.def))
(/usr/share/texlive/texmf-dist/tex/latex/bxjscls/bxjsja-standard.def
(/usr/share/texlive/texmf-dist/tex/latex/bxjscls/bxjsja-minimal.def)
(/usr/share/texlive/texmf-dist/tex/latex/zxjatype/zxjatype.sty
(/usr/share/texlive/texmf-dist/tex/generic/iftex/ifxetex.sty)
(/usr/share/texlive/texmf-dist/tex/xelatex/xecjk/xeCJK.sty
(/usr/share/texlive/texmf-dist/tex/latex/l3kernel/expl3.sty
(/usr/share/texlive/texmf-dist/tex/latex/l3backend/l3backend-xetex.def))
(/usr/share/texlive/texmf-dist/tex/latex/ctex/ctexhook.sty)
(/usr/share/texlive/texmf-dist/tex/latex/l3packages/xtemplate/xtemplate.sty)
(/usr/share/texlive/texmf-dist/tex/latex/fontspec/fontspec.sty
(/usr/share/texlive/texmf-dist/tex/latex/l3packages/xparse/xparse.sty)
(/usr/share/texlive/texmf-dist/tex/latex/fontspec/fontspec-xetex.sty
(/usr/share/texlive/texmf-dist/tex/latex/base/fontenc.sty)
(/usr/share/texlive/texmf-dist/tex/latex/fontspec/fontspec.cfg)))
(/usr/share/texlive/texmf-dist/tex/xelatex/xecjk/xeCJK.cfg)))))
No file xelatex-test.aux.
(/usr/share/texlive/texmf-dist/tex/latex/base/ts1cmr.fd)
*geometry* driver: auto-detecting
*geometry* detected driver: xetex
[1] (./xelatex-test.aux) )
Output written on xelatex-test.pdf (1 page).
Transcript written on xelatex-test.log.
This is pdfTeX, Version 3.141592653-2.6-1.40.24 (TeX Live 2022/Debian) (preloaded format=pdflatex)
 restricted \write18 enabled.
entering extended mode
(./pdflatex-test.tex
LaTeX2e <2022-11-01> patch level 1
L3 programming layer <2023-01-16>
(/usr/share/texlive/texmf-dist/tex/latex/base/article.cls
Document Class: article 2022/07/02 v1.4n Standard LaTeX document class
(/usr/share/texlive/texmf-dist/tex/latex/base/size10.clo))
(/usr/share/texlive/texmf-dist/tex/latex/l3backend/l3backend-pdftex.def)
No file pdflatex-test.aux.
[1{/var/lib/texmf/fonts/map/pdftex/updmap/pdftex.map}] (./pdflatex-test.aux) )<
/usr/share/texlive/texmf-dist/fonts/type1/public/amsfonts/cm/cmmi10.pfb></usr/s
hare/texlive/texmf-dist/fonts/type1/public/amsfonts/cm/cmmi7.pfb></usr/share/te
xlive/texmf-dist/fonts/type1/public/amsfonts/cm/cmr10.pfb>
Output written on pdflatex-test.pdf (1 page, 30138 bytes).
Transcript written on pdflatex-test.log.
This is e-upTeX, Version 3.141592653-p4.0.0-u1.28-220214-2.6 (utf8.uptex) (TeX Live 2022/Debian) (preloaded format=uplatex)
 restricted \write18 enabled.
entering extended mode
(./uplatex-test.tex
pLaTeX2e <2021-11-15u04>+1 (based on LaTeX2e <2022-11-01> patch level 1)
L3 programming layer <2023-01-16>
(/usr/share/texlive/texmf-dist/tex/platex/jsclasses/jsarticle.cls
Document Class: jsarticle 2022/09/13 jsclasses (okumura, texjporg)
(/usr/share/texlive/texmf-dist/tex/platex/jsclasses/jslogo.sty))
(/usr/share/texlive/texmf-dist/tex/latex/l3backend/l3backend-dvipdfmx.def)
No file uplatex-test.aux.
[1] (./uplatex-test.aux) )
Output written on uplatex-test.dvi (1 page, 420 bytes).
Transcript written on uplatex-test.log.
uplatex-test.dvi -> uplatex-test.pdf
[1]
5268 bytes written
This is e-pTeX, Version 3.141592653-p4.0.0-220214-2.6 (utf8.euc) (TeX Live 2022/Debian) (preloaded format=platex)
 restricted \write18 enabled.
entering extended mode
(./platex-test.tex
pLaTeX2e <2021-11-15>+1 (based on LaTeX2e <2022-11-01> patch level 1)
L3 programming layer <2023-01-16>
(/usr/share/texlive/texmf-dist/tex/platex/jsclasses/jsarticle.cls
Document Class: jsarticle 2022/09/13 jsclasses (okumura, texjporg)
(/usr/share/texlive/texmf-dist/tex/platex/jsclasses/jslogo.sty))
(/usr/share/texlive/texmf-dist/tex/latex/l3backend/l3backend-dvipdfmx.def)
No file platex-test.aux.
[1] (./platex-test.aux) )
Output written on platex-test.dvi (1 page, 408 bytes).
Transcript written on platex-test.log.
platex-test.dvi -> platex-test.pdf
[1]
5168 bytes written
-rw-r--r-- 1 mat mat  28K Aug 14 22:56 lualatex-test.pdf
-rw-r--r-- 1 mat mat  30K Aug 14 22:56 pdflatex-test.pdf
-rw-r--r-- 1 mat mat 5.1K Aug 14 22:56 platex-test.pdf
-rw-r--r-- 1 mat mat 5.2K Aug 14 22:56 uplatex-test.pdf
-rw-r--r-- 1 mat mat 8.4K Aug 14 22:56 xelatex-test.pdf
lualatex-test.pdf: PDF document, version 1.5
pdflatex-test.pdf: PDF document, version 1.5
platex-test.pdf:   PDF document, version 1.5 (zip deflate encoded)
uplatex-test.pdf:  PDF document, version 1.5 (zip deflate encoded)
xelatex-test.pdf:  PDF document, version 1.5 (zip deflate encoded)
\end{verbatim}
\end{comment}

\section{Python版組API}

\subsection{仮想環境}

Raspberry Pi上で実行する。

%command4.sh

\begin{command}
sudo apt install python3-venv python3-pip openssl
mkdir -p /home/mat/texnote-server
cd /home/mat/texnote-server
python3 -m venv .venv
source .venv/bin/activate
python -m pip install --upgrade pip
python -m pip install 'fastapi==0.141.1' 'uvicorn[standard]==0.52.1'
python --version
python -m pip show fastapi | grep Version
python -m pip show uvicorn | grep Version
which python
which uvicorn
deactivate
\end{command}

\begin{comment}
command4.sh の実行結果

\begin{verbatim}
Reading package lists... Done
Building dependency tree... Done
Reading state information... Done
python3-venv is already the newest version (3.11.2-1+b1).
python3-pip is already the newest version (23.0.1+dfsg-1+rpt1).
openssl is already the newest version (3.0.20-1~deb12u2+rpt1).
The following packages were automatically installed and are no longer required:
  avahi-utils gir1.2-handy-1 gir1.2-packagekitglib-1.0 gir1.2-polkit-1.0
  libbasicusageenvironment1 libc++1-16 libc++abi1-16 libcamera0.4
  libgroupsock8 liblivemedia77 libunwind-16
  linux-headers-6.6.62+rpt-common-rpi linux-headers-6.6.62+rpt-rpi-2712
  linux-headers-6.6.62+rpt-rpi-v8 linux-headers-6.6.74+rpt-common-rpi
  linux-headers-6.6.74+rpt-rpi-2712 linux-headers-6.6.74+rpt-rpi-v8
  linux-image-6.6.62+rpt-rpi-2712 linux-image-6.6.62+rpt-rpi-v8
  linux-image-6.6.74+rpt-rpi-2712 linux-image-6.6.74+rpt-rpi-v8
  linux-kbuild-6.6.62+rpt linux-kbuild-6.6.74+rpt python3-v4l2
Use 'sudo apt autoremove' to remove them.
0 upgraded, 0 newly installed, 0 to remove and 0 not upgraded.
command4.sh: 5: source: not found
error: externally-managed-environment

× This environment is externally managed
╰─> To install Python packages system-wide, try apt install
    python3-xyz, where xyz is the package you are trying to
    install.

    If you wish to install a non-Debian-packaged Python package,
    create a virtual environment using python3 -m venv path/to/venv.
    Then use path/to/venv/bin/python and path/to/venv/bin/pip. Make
    sure you have python3-full installed.

    For more information visit http://rptl.io/venv

note: If you believe this is a mistake, please contact your Python installation or OS distribution provider. You can override this, at the risk of breaking your Python installation or OS, by passing --break-system-packages.
hint: See PEP 668 for the detailed specification.
error: externally-managed-environment

× This environment is externally managed
╰─> To install Python packages system-wide, try apt install
    python3-xyz, where xyz is the package you are trying to
    install.

    If you wish to install a non-Debian-packaged Python package,
    create a virtual environment using python3 -m venv path/to/venv.
    Then use path/to/venv/bin/python and path/to/venv/bin/pip. Make
    sure you have python3-full installed.

    For more information visit http://rptl.io/venv

note: If you believe this is a mistake, please contact your Python installation or OS distribution provider. You can override this, at the risk of breaking your Python installation or OS, by passing --break-system-packages.
hint: See PEP 668 for the detailed specification.
Python 3.11.2
WARNING: Package(s) not found: fastapi
WARNING: Package(s) not found: uvicorn
/usr/bin/python
command4.sh: 13: deactivate: not found
mat@raspberrypi:~ $ source .venv/bin/activate
-bash: .venv/bin/activate: No such file or directory
mat@raspberrypi:~ $ ls ./.venv
ls: cannot access './.venv': No such file or directory
mat@raspberrypi:~ $ cd texnote-server/
mat@raspberrypi:~/texnote-server $ ls
mat@raspberrypi:~/texnote-server $ ls ./.venv
bin  include  lib  lib64  pyvenv.cfg
mat@raspberrypi:~/texnote-server $ source .venv/bin/activate
(.venv) mat@raspberrypi:~/texnote-server $ python -m pip install --upgrade pip
Looking in indexes: https://pypi.org/simple, https://www.piwheels.org/simple
Requirement already satisfied: pip in ./.venv/lib/python3.11/site-packages (23.0.1)
Collecting pip
  Obtaining dependency information for pip from https://www.piwheels.org/simple/pip/pip-26.2.1-py3-none-any.whl.metadata
  Downloading https://www.piwheels.org/simple/pip/pip-26.2.1-py3-none-any.whl.metadata (4.6 kB)
Downloading https://www.piwheels.org/simple/pip/pip-26.2.1-py3-none-any.whl (1.8 MB)
   ━━━━━━━━━━━━━━━━━━━━━━━━━━━━━━━━━━━━━━━━ 1.8/1.8 MB 736.2 kB/s eta 0:00:00
Using cached https://www.piwheels.org/simple/pip/pip-26.2.1-py3-none-any.whl (1.8 MB)
Installing collected packages: pip
  Attempting uninstall: pip
    Found existing installation: pip 23.0.1
    Uninstalling pip-23.0.1:
      Successfully uninstalled pip-23.0.1
Successfully installed pip-26.2.1
(.venv) mat@raspberrypi:~/texnote-server $ python -m pip install 'fastapi==0.141.1' 'uvicorn[standard]==0.52.1'
Looking in indexes: https://pypi.org/simple, https://www.piwheels.org/simple
Collecting fastapi==0.141.1
  Downloading fastapi-0.141.1-py3-none-any.whl.metadata (27 kB)
Collecting uvicorn==0.52.1 (from uvicorn[standard]==0.52.1)
  Downloading uvicorn-0.52.1-py3-none-any.whl.metadata (6.6 kB)
Collecting starlette>=0.46.0 (from fastapi==0.141.1)
  Downloading starlette-1.6.0-py3-none-any.whl.metadata (6.4 kB)
Collecting pydantic>=2.9.0 (from fastapi==0.141.1)
  Downloading pydantic-2.13.4-py3-none-any.whl.metadata (109 kB)
Collecting typing-extensions>=4.8.0 (from fastapi==0.141.1)
  Downloading typing_extensions-4.16.0-py3-none-any.whl.metadata (3.3 kB)
Collecting typing-inspection>=0.4.2 (from fastapi==0.141.1)
  Downloading typing_inspection-0.4.4-py3-none-any.whl.metadata (2.6 kB)
Collecting annotated-doc>=0.0.2 (from fastapi==0.141.1)
  Downloading annotated_doc-0.0.5-py3-none-any.whl.metadata (6.5 kB)
Collecting click>=7.0 (from uvicorn==0.52.1->uvicorn[standard]==0.52.1)
  Downloading click-8.4.2-py3-none-any.whl.metadata (2.6 kB)
Collecting h11>=0.8 (from uvicorn==0.52.1->uvicorn[standard]==0.52.1)
  Downloading h11-0.16.0-py3-none-any.whl.metadata (8.3 kB)
Collecting httptools>=0.8.0 (from uvicorn[standard]==0.52.1)
  Downloading httptools-0.8.0-cp311-cp311-manylinux2014_aarch64.manylinux_2_17_aarch64.manylinux_2_28_aarch64.whl.metadata (3.5 kB)
Collecting python-dotenv>=0.13 (from uvicorn[standard]==0.52.1)
  Downloading python_dotenv-1.2.2-py3-none-any.whl.metadata (27 kB)
Collecting pyyaml>=5.1 (from uvicorn[standard]==0.52.1)
  Downloading pyyaml-6.0.3-cp311-cp311-manylinux2014_aarch64.manylinux_2_17_aarch64.manylinux_2_28_aarch64.whl.metadata (2.4 kB)
Collecting uvloop>=0.15.1 (from uvicorn[standard]==0.52.1)
  Downloading uvloop-0.22.1-cp311-cp311-manylinux2014_aarch64.manylinux_2_17_aarch64.manylinux_2_28_aarch64.whl.metadata (4.9 kB)
Collecting watchfiles>=0.20 (from uvicorn[standard]==0.52.1)
  Downloading watchfiles-1.2.0-cp311-cp311-manylinux_2_17_aarch64.manylinux2014_aarch64.whl.metadata (4.9 kB)
Collecting websockets>=13.0 (from uvicorn[standard]==0.52.1)
  Downloading websockets-17.0.1-cp311-cp311-manylinux2014_aarch64.manylinux_2_17_aarch64.manylinux_2_28_aarch64.whl.metadata (6.3 kB)
Collecting annotated-types>=0.6.0 (from pydantic>=2.9.0->fastapi==0.141.1)
  Downloading annotated_types-0.8.0-py3-none-any.whl.metadata (15 kB)
Collecting pydantic-core==2.46.4 (from pydantic>=2.9.0->fastapi==0.141.1)
  Downloading pydantic_core-2.46.4-cp311-cp311-manylinux_2_17_aarch64.manylinux2014_aarch64.whl.metadata (6.6 kB)
Collecting anyio<5,>=3.6.2 (from starlette>=0.46.0->fastapi==0.141.1)
  Downloading anyio-4.14.2-py3-none-any.whl.metadata (4.6 kB)
Collecting idna>=2.8 (from anyio<5,>=3.6.2->starlette>=0.46.0->fastapi==0.141.1)
  Downloading idna-3.18-py3-none-any.whl.metadata (6.1 kB)
Downloading fastapi-0.141.1-py3-none-any.whl (131 kB)
Downloading uvicorn-0.52.1-py3-none-any.whl (79 kB)
Downloading annotated_doc-0.0.5-py3-none-any.whl (5.3 kB)
Downloading click-8.4.2-py3-none-any.whl (119 kB)
Downloading h11-0.16.0-py3-none-any.whl (37 kB)
Downloading httptools-0.8.0-cp311-cp311-manylinux2014_aarch64.manylinux_2_17_aarch64.manylinux_2_28_aarch64.whl (464 kB)
Downloading pydantic-2.13.4-py3-none-any.whl (472 kB)
Downloading pydantic_core-2.46.4-cp311-cp311-manylinux_2_17_aarch64.manylinux2014_aarch64.whl (2.0 MB)
   ━━━━━━━━━━━━━━━━━━━━━━━━━━━━━━━━━━━━━━━━ 2.0/2.0 MB 3.8 MB/s  0:00:00
Downloading annotated_types-0.8.0-py3-none-any.whl (13 kB)
Downloading python_dotenv-1.2.2-py3-none-any.whl (22 kB)
Downloading pyyaml-6.0.3-cp311-cp311-manylinux2014_aarch64.manylinux_2_17_aarch64.manylinux_2_28_aarch64.whl (775 kB)
   ━━━━━━━━━━━━━━━━━━━━━━━━━━━━━━━━━━━━━━━━ 775.6/775.6 kB 4.0 MB/s  0:00:00
Downloading starlette-1.6.0-py3-none-any.whl (75 kB)
Downloading anyio-4.14.2-py3-none-any.whl (124 kB)
Downloading idna-3.18-py3-none-any.whl (65 kB)
Downloading typing_extensions-4.16.0-py3-none-any.whl (45 kB)
Downloading typing_inspection-0.4.4-py3-none-any.whl (14 kB)
Downloading uvloop-0.22.1-cp311-cp311-manylinux2014_aarch64.manylinux_2_17_aarch64.manylinux_2_28_aarch64.whl (3.8 MB)
   ━━━━━━━━━━━━━━━━━━━━━━━━━━━━━━━━━━━━━━━━ 3.8/3.8 MB 3.8 MB/s  0:00:00
Downloading watchfiles-1.2.0-cp311-cp311-manylinux_2_17_aarch64.manylinux2014_aarch64.whl (456 kB)
Downloading websockets-17.0.1-cp311-cp311-manylinux2014_aarch64.manylinux_2_17_aarch64.manylinux_2_28_aarch64.whl (221 kB)
Installing collected packages: websockets, uvloop, typing-extensions, pyyaml, python-dotenv, idna, httptools, h11, click, annotated-types, annotated-doc, uvicorn, typing-inspection, pydantic-core, anyio, watchfiles, starlette, pydantic, fastapi
Successfully installed annotated-doc-0.0.5 annotated-types-0.8.0 anyio-4.14.2 click-8.4.2 fastapi-0.141.1 h11-0.16.0 httptools-0.8.0 idna-3.18 pydantic-2.13.4 pydantic-core-2.46.4 python-dotenv-1.2.2 pyyaml-6.0.3 starlette-1.6.0 typing-extensions-4.16.0 typing-inspection-0.4.4 uvicorn-0.52.1 uvloop-0.22.1 watchfiles-1.2.0 websockets-17.0.1
(.venv) mat@raspberrypi:~/texnote-server $ python --version
Python 3.11.2
(.venv) mat@raspberrypi:~/texnote-server $ python -m pip show fastapi | grep Version
Version: 0.141.1
(.venv) mat@raspberrypi:~/texnote-server $ python -m pip show uvicorn | grep Version
Version: 0.52.1
(.venv) mat@raspberrypi:~/texnote-server $ which python
/home/mat/texnote-server/.venv/bin/python
(.venv) mat@raspberrypi:~/texnote-server $ which uvicorn
/home/mat/texnote-server/.venv/bin/uvicorn
(.venv) mat@raspberrypi:~/texnote-server $ deactivate
mat@raspberrypi:~/texnote-server $
\end{verbatim}
\end{comment}

再現用\texttt{requirements.txt}は次の内容である。
\begin{command}
fastapi==0.141.1
uvicorn[standard]==0.52.1
\end{command}

\subsection{Macからプログラムを転送}

Mac側のリポジトリ位置が異なる場合はパスを変更する。
\begin{command}
scp \
  /Users/mat/Documents/Git/TeXNote/Server/P2/progs/app.py \
  /Users/mat/Documents/Git/TeXNote/Server/P2/progs/requirements.txt \
  mat@192.168.3.11:/home/mat/texnote-server/
\end{command}
Raspberry Piで依存関係をそろえる。
\begin{command}
cd /home/mat/texnote-server
source .venv/bin/activate
python -m pip install -r requirements.txt
\end{command}

\begin{comment}
app.py と requirements.txt をコピー後

\begin{verbatim}
mat@raspberrypi:~/texnote-server $ ls
app.py  requirements.txt
mat@raspberrypi:~/texnote-server $ cat requirements.txt
fastapi==0.141.1
uvicorn[standard]==0.52.1
mat@raspberrypi:~/texnote-server $ source .venv/bin/activate
(.venv) mat@raspberrypi:~/texnote-server $ python -m pip install -r requirements.txt
Looking in indexes: https://pypi.org/simple, https://www.piwheels.org/simple
Requirement already satisfied: fastapi==0.141.1 in ./.venv/lib/python3.11/site-packages (from -r requirements.txt (line 1)) (0.141.1)
Requirement already satisfied: uvicorn==0.52.1 in ./.venv/lib/python3.11/site-packages (from uvicorn[standard]==0.52.1->-r requirements.txt (line 2)) (0.52.1)
Requirement already satisfied: starlette>=0.46.0 in ./.venv/lib/python3.11/site-packages (from fastapi==0.141.1->-r requirements.txt (line 1)) (1.6.0)
Requirement already satisfied: pydantic>=2.9.0 in ./.venv/lib/python3.11/site-packages (from fastapi==0.141.1->-r requirements.txt (line 1)) (2.13.4)
Requirement already satisfied: typing-extensions>=4.8.0 in ./.venv/lib/python3.11/site-packages (from fastapi==0.141.1->-r requirements.txt (line 1)) (4.16.0)
Requirement already satisfied: typing-inspection>=0.4.2 in ./.venv/lib/python3.11/site-packages (from fastapi==0.141.1->-r requirements.txt (line 1)) (0.4.4)
Requirement already satisfied: annotated-doc>=0.0.2 in ./.venv/lib/python3.11/site-packages (from fastapi==0.141.1->-r requirements.txt (line 1)) (0.0.5)
Requirement already satisfied: click>=7.0 in ./.venv/lib/python3.11/site-packages (from uvicorn==0.52.1->uvicorn[standard]==0.52.1->-r requirements.txt (line 2)) (8.4.2)
Requirement already satisfied: h11>=0.8 in ./.venv/lib/python3.11/site-packages (from uvicorn==0.52.1->uvicorn[standard]==0.52.1->-r requirements.txt (line 2)) (0.16.0)
Requirement already satisfied: httptools>=0.8.0 in ./.venv/lib/python3.11/site-packages (from uvicorn[standard]==0.52.1->-r requirements.txt (line 2)) (0.8.0)
Requirement already satisfied: python-dotenv>=0.13 in ./.venv/lib/python3.11/site-packages (from uvicorn[standard]==0.52.1->-r requirements.txt (line 2)) (1.2.2)
Requirement already satisfied: pyyaml>=5.1 in ./.venv/lib/python3.11/site-packages (from uvicorn[standard]==0.52.1->-r requirements.txt (line 2)) (6.0.3)
Requirement already satisfied: uvloop>=0.15.1 in ./.venv/lib/python3.11/site-packages (from uvicorn[standard]==0.52.1->-r requirements.txt (line 2)) (0.22.1)
Requirement already satisfied: watchfiles>=0.20 in ./.venv/lib/python3.11/site-packages (from uvicorn[standard]==0.52.1->-r requirements.txt (line 2)) (1.2.0)
Requirement already satisfied: websockets>=13.0 in ./.venv/lib/python3.11/site-packages (from uvicorn[standard]==0.52.1->-r requirements.txt (line 2)) (17.0.1)
Requirement already satisfied: annotated-types>=0.6.0 in ./.venv/lib/python3.11/site-packages (from pydantic>=2.9.0->fastapi==0.141.1->-r requirements.txt (line 1)) (0.8.0)
Requirement already satisfied: pydantic-core==2.46.4 in ./.venv/lib/python3.11/site-packages (from pydantic>=2.9.0->fastapi==0.141.1->-r requirements.txt (line 1)) (2.46.4)
Requirement already satisfied: anyio<5,>=3.6.2 in ./.venv/lib/python3.11/site-packages (from starlette>=0.46.0->fastapi==0.141.1->-r requirements.txt (line 1)) (4.14.2)
Requirement already satisfied: idna>=2.8 in ./.venv/lib/python3.11/site-packages (from anyio<5,>=3.6.2->starlette>=0.46.0->fastapi==0.141.1->-r requirements.txt (line 1)) (3.18)
(.venv) mat@raspberrypi:~/texnote-server $ deactivate
mat@raspberrypi:~/texnote-server $
\end{verbatim}
\end{comment}

\section{APIトークン}

Raspberry Pi上で作成する。トークンはGitへ登録しない。

%command6.sh

\begin{command}
cd /home/mat/texnote-server
umask 077
openssl rand -hex 32 > .api-token
chmod 600 .api-token
printf 'TEXNOTE_API_TOKEN=' > server.env
cat .api-token >> server.env
chmod 600 server.env
ls -l .api-token server.env
\end{command}

\begin{comment}
\begin{verbatim}
mat@raspberrypi:~ $ sh command6.sh
-rw------- 1 mat mat 65 Aug 14 23:11 .api-token
-rw------- 1 mat mat 83 Aug 14 23:11 server.env
mat@raspberrypi:~ $
\end{verbatim}
\end{comment}

\section{手動起動と疎通確認}

Raspberry Piで起動する。

\begin{command}
cd /home/mat/texnote-server
source .venv/bin/activate
set -a
source server.env
set +a
uvicorn app:app --host 192.168.3.11 --port 8000
\end{command}

Macの別ターミナルで確認する。
\begin{command}
curl --fail-with-body http://192.168.3.11:8000/health
\end{command}
期待される応答例：
\begin{command}
{"status":"ok","engines":["lualatex","xelatex","pdflatex","uplatex","platex"]}
\end{command}
認証も確認し、Mac上の一時変数を最後に消去する。
\begin{command}
TEXNOTE_TOKEN="$(ssh mat@192.168.3.11 \
  'cat /home/mat/texnote-server/.api-token')"
curl --fail-with-body \
  http://192.168.3.11:8000/v1/auth-check \
  -H "Authorization: Bearer $TEXNOTE_TOKEN"
unset TEXNOTE_TOKEN
\end{command}
成功時は\texttt{\{"status":"ok"\}}が返る。手動サーバーはPi側でCtrl-Cを押して終了する。

\section{APIリクエストと添付ファイルの制限}

P2は過大な入力によるメモリ・CPU・一時領域の消費を抑えるため、現行の
\texttt{app.py}で次の上限を設定している。

\begin{longtable}{@{}p{58mm}p{38mm}p{54mm}@{}}
\toprule
対象 & 現行上限 & 超過時の扱い\\
\midrule
HTTPリクエスト全体 & 32 MiB & HTTP 413\\
TeXソース & 2,000,000文字 & 入力検証エラー\\
添付1ファイル & 10 MiB & HTTP 413\\
添付ファイル合計 & 20 MiB & HTTP 413\\
\texttt{pictures}の個数 & 50個 & 入力検証エラー\\
\texttt{files}の個数 & 50個 & 入力検証エラー\\
相対パス & 1,024文字 & 入力検証エラー\\
P2ログ & 200,000文字 & 末尾を応答に使用\\
\bottomrule
\end{longtable}

HTTPリクエスト全体の上限は次の定数で33,554,432 bytesに設定する。

\begin{command}
MAX_REQUEST_BYTES = 32 * 1024 * 1024
\end{command}

画像と添付ファイルはJSONへBase64文字列として格納される。Base64は元データのおよそ
4/3倍になるため、復号後20 MiBの添付は約26.7 MiBになる。これにTeXソース、相対パス、
JSON構造を加えても収まるよう、HTTP全体を32 MiBとしている。

\texttt{Content-Length}が32 MiBを超える場合は、本文を解析する前にHTTP 413を返す。

\begin{command}
{"message":"リクエストの上限32 MiBを超えています。"}
\end{command}

P1にも同じ32 MiB制限を設定し、P1で受理された要求がP2の入口で容量超過にならないよう
上限を統一する。公開時にリバースプロキシを置く場合も、そこで同じ上限を設定する。

\section{systemdによる自動起動}

Raspberry Pi上で作成する。
\begin{command}
sudo nano /etc/systemd/system/texnote-server.service
\end{command}
内容は次のとおり。\texttt{ExecStart}は1行で記述する。
\begin{command}
[Unit]
Description=TeXNote Typesetting Server
After=network-online.target
Wants=network-online.target

[Service]
Type=simple
User=mat
Group=mat
WorkingDirectory=/home/mat/texnote-server
EnvironmentFile=/home/mat/texnote-server/server.env
ExecStart=/home/mat/texnote-server/.venv/bin/uvicorn app:app --host 192.168.3.11 --port 8000
Restart=on-failure
RestartSec=3
NoNewPrivileges=true
PrivateTmp=true

[Install]
WantedBy=multi-user.target
\end{command}
読み込み、自動起動を有効化する。
\begin{command}
sudo systemctl daemon-reload
sudo systemctl enable --now texnote-server
systemctl status texnote-server --no-pager
systemctl is-enabled texnote-server
systemctl is-active texnote-server
journalctl -u texnote-server -n 100 --no-pager
\end{command}
再起動試験をする。
\begin{command}
sudo reboot
\end{command}
Piが戻った後、Macで確認する。
\begin{command}
ping -c 3 192.168.3.11
curl --connect-timeout 5 http://192.168.3.11:8000/health
ssh mat@192.168.3.11 \
  'systemctl is-enabled texnote-server && systemctl is-active texnote-server'
\end{command}

\section{iPad版・iPhone版の設定}
アプリへ次を登録する。
\begin{itemize}
\item サーバーURL：\texttt{http://192.168.3.11:8000}
\item APIトークン：Piの\texttt{/home/mat/texnote-server/.api-token}の内容
\end{itemize}
現段階は同一Wi-Fi内だけで使用する。ルーターの8000番ポートをインターネットへ
転送しない。外出先利用はTailscaleやWireGuard等のVPNを導入してから行う。

\section{プログラムの更新}
\subsection{systemd版の更新}

アプリと版組サーバーのAPIを変更した場合は、両者の対応版を同時期に更新する。
特に添付リソースのキーを\texttt{fileName}から\texttt{relativePath}へ変更した版では、
新旧のアプリとサーバーを混在させると版組要求がHTTP 422で失敗する。

まずMacから新しい\texttt{app.py}を一時名で転送する。
\begin{command}
scp \
  /Users/mat/Documents/Git/TeXNote/Server/P2/progs/app.py \
  mat@192.168.3.11:/home/mat/texnote-server/app.py.new
\end{command}
Pi上で現行版をバックアップし、差替え前にPython構文検査を行う。
\begin{command}
ssh mat@192.168.3.11
cd /home/mat/texnote-server
cp -p app.py app.py.backup
.venv/bin/python -m py_compile app.py.new
mv app.py.new app.py
\end{command}
サービスを再起動し、状態、ログ、HTTP応答を確認する。
\begin{command}
sudo systemctl restart texnote-server
sudo systemctl status texnote-server --no-pager
sudo journalctl -u texnote-server -n 100 --no-pager
curl --fail-with-body http://192.168.3.11:8000/health
\end{command}
異常があればバックアップへ戻す。
\begin{command}
cd /home/mat/texnote-server
cp -p app.py.backup app.py
sudo systemctl restart texnote-server
sudo systemctl status texnote-server --no-pager
\end{command}
サーバーの\texttt{/health}確認後に対応するiPad・iPhone版アプリを更新し、
サブフォルダ内の画像または添付ファイルを含むCardで版組、保存、再読込を確認する。

\section{画像・ファイル・独自スタイル}
文書固有の\texttt{.sty}、画像、データはTeXNoteの「画像・ファイル」から送る。
新しいAPIでは各リソースがノートルート基準の安全な相対パスを持ち、
サーバー上の一時作業ディレクトリにもその階層が再現される。
例えば\texttt{files/}以下のリソースは次のように参照する。
\begin{command}
\usepackage{files/my-style}
\includegraphics{files/example.png}
\input{files/part.tex}
\end{command}
全文書で共通利用するスタイルをPiのローカルTeXツリーへ置く例：
\begin{command}
mkdir -p ~/texmf/tex/latex/texnote
cp my-style.sty ~/texmf/tex/latex/texnote/
mktexlsr ~/texmf
kpsewhich my-style.sty
\end{command}

\section{pdfLaTeX/CJK用Type 1フォント一式}
\subsection{症状と事前確認}
次のようなエラーは、IPAexのTrueTypeフォントだけでなく、
pdfLaTeX/CJKが使うTFM/PFB/mapが不足していることを示す。
\begin{command}
I can't find file 'ipxm-r-u61'
Font ... =ipxm-r-u61 ... not loadable:
Metric (TFM) file not found.
Fatal error occurred, no output PDF file produced!
\end{command}
\begin{command}
kpsewhich ipaexm.ttf
kpsewhich ipaexg.ttf
kpsewhich ipxm-r-u61.tfm
\end{command}

\subsection{導入}
IPAexフォントを確認・導入する。
\begin{command}
sudo apt update
sudo apt install fonts-ipaexfont-mincho fonts-ipaexfont-gothic
\end{command}
DebianではType 1フォント一式を\texttt{texlive-fonts-extra}で導入する。
容量が大きいため、先に空き容量を確認する。
\begin{command}
df -h /
sudo apt install texlive-fonts-extra
sudo mktexlsr
sudo updmap-sys
\end{command}
導入結果を確認する。
\begin{command}
kpsewhich ipaexm.ttf
kpsewhich ipaexg.ttf
kpsewhich ipxm-r-u61.tfm
kpsewhich ipxm-r-u61.pfb
kpsewhich ipaex-type1.map
\end{command}
各コマンドでパスが表示されればTeX Liveから発見できている。
\begin{command}
sudo systemctl restart texnote-server
systemctl is-active texnote-server
\end{command}

\section{日常の確認}
\begin{command}
systemctl status texnote-server --no-pager
journalctl -u texnote-server -n 100 --no-pager
journalctl -u texnote-server -f
curl --fail-with-body http://192.168.3.11:8000/health
df -h /
free -h
kpsewhich ファイル名.sty
\end{command}

\section{完了チェックリスト}
\begin{itemize}
\item[$\square$] OS、aarch64、容量、メモリ、IPを確認した。
\item[$\square$] IPアドレスをDHCP予約した。
\item[$\square$] 5種類のTeXエンジンがPDFを生成した。
\item[$\square$] 仮想環境へFastAPIとUvicornを導入した。
\item[$\square$] APIトークンの権限を600にした。
\item[$\square$] \texttt{/health}と\texttt{/v1/auth-check}が成功した。
\item[$\square$] systemdがenabledかつactiveになった。
\item[$\square$] Pi再起動後も自動復旧した。
\item[$\square$] iPad版・iPhone版の双方から版組できた。
\item[$\square$] Type 1のTFM/PFB/mapを確認した。
\item[$\square$] TCP 8000をインターネットへ公開していない。
\end{itemize}

\newpage
\section{Docker版への移行}

ここまでは、Raspberry Pi本体へTeX LiveとPython仮想環境を導入し、
systemdからUvicornを直接起動する構成を説明した。
本節では、同じ版組APIとTeX環境をDockerイメージへまとめる。

Docker版はホスト側のPython仮想環境やTeX Liveに依存せず、
\texttt{Dockerfile}、\texttt{compose.yaml}、\texttt{app.py}、
\texttt{requirements.txt}から再構築できる。
一方、TeX Liveと\texttt{texlive-fonts-extra}を含むため、イメージは大きくなり、
初回ビルドには長い時間と数GB以上の空き容量を要する。

\subsection{移行方針と注意事項}

本手順では次の構成にする。
\begin{itemize}
\item ベースイメージはarm64対応の\texttt{python:3.11-slim-bookworm}を使う。
\item TeX Live、日本語環境、IPAex、Type 1フォントをイメージへ組み込む。
\item APIトークンはイメージへコピーせず、起動時に\texttt{server.env}から渡す。
\item コンテナ内ではrootではなく専用ユーザー\texttt{texnote}で実行する。
\item コンテナは読み取り専用とし、版組用一時ファイルだけを\texttt{/tmp}へ置く。
\item ホストの\texttt{192.168.3.11:8000}だけをコンテナへ転送する。
\item 現在のsystemd版は、Docker版のビルドと検査が終わるまで削除しない。
\end{itemize}

systemd版とDocker版は同じTCP 8000を同時に使用できない。
Docker版を起動する直前にsystemd版を停止し、問題があれば直ちに戻せるようにする。
また、Dockerが公開するポートはufwの規則を迂回する場合がある。
本例のようにLAN側IPアドレスへ明示的にbindし、ルーターでは8000番を転送しない。

\subsection{Docker EngineとComposeの導入}

以下はRaspberry Pi上で実行する。Docker公式のDebian用APTリポジトリを使う。
Debian 12とarm64はDocker Engineのサポート対象である。

既存の競合パッケージがないか確認する。
\begin{command}
dpkg -l docker.io docker-compose docker-doc docker-buildx \
  podman-docker containerd runc 2>/dev/null || true
\end{command}
既に別のDocker環境を運用している場合は、ここで削除せず構成を確認する。
未導入の環境では、公式鍵とリポジトリを登録する。
\begin{command}
sudo apt update
sudo apt install ca-certificates curl
sudo install -m 0755 -d /etc/apt/keyrings
sudo curl -fsSL https://download.docker.com/linux/debian/gpg \
  -o /etc/apt/keyrings/docker.asc
sudo chmod a+r /etc/apt/keyrings/docker.asc

sudo tee /etc/apt/sources.list.d/docker.sources >/dev/null <<EOF
Types: deb
URIs: https://download.docker.com/linux/debian
Suites: bookworm
Components: stable
Architectures: arm64
Signed-By: /etc/apt/keyrings/docker.asc
EOF

sudo apt update
sudo apt install \
  docker-ce \
  docker-ce-cli \
  containerd.io \
  docker-buildx-plugin \
  docker-compose-plugin
\end{command}
導入結果を確認する。
\begin{command}
sudo systemctl is-enabled docker
sudo systemctl is-active docker
sudo docker version
sudo docker compose version
sudo docker run --rm hello-world
\end{command}

本書ではすべて\texttt{sudo docker}で実行する。
ユーザーを\texttt{docker}グループへ追加すると\texttt{sudo}なしで操作できるが、
そのユーザーには実質的にroot相当の権限が与えられるため、ここでは行わない。

\subsection{Docker用ディレクトリの準備}

既存のsystemd版と混同しないよう、Pi上に別ディレクトリを作る。
\begin{command}
mkdir -p /home/mat/texnote-server-docker
chmod 700 /home/mat/texnote-server-docker
\end{command}

Mac側から現在のサーバープログラムを転送する。
\begin{command}
scp \
  /Users/mat/Documents/Git/TeXNote/Server/P2/progs/app.py \
  /Users/mat/Documents/Git/TeXNote/Server/P2/progs/requirements.txt \
  mat@192.168.3.11:/home/mat/texnote-server-docker/
\end{command}

APIトークンは既存の値を引き継げる。Pi上でコピーし、所有者以外から読めないようにする。
\begin{command}
sudo install -o mat -g mat -m 600 \
  /home/mat/texnote-server/server.env \
  /home/mat/texnote-server-docker/server.env
\end{command}
新しいトークンへ交換する場合は次のように作る。
この場合、iPad版とiPhone版の登録値も更新する。
\begin{command}
cd /home/mat/texnote-server-docker
umask 077
printf 'TEXNOTE_API_TOKEN=' > server.env
openssl rand -hex 32 >> server.env
chmod 600 server.env
\end{command}

\subsection{Dockerfileの作成}

Pi上の\texttt{/home/mat/texnote-server-docker/Dockerfile}を次の内容で作成する。
\begin{command}
cd /home/mat/texnote-server-docker
nano Dockerfile
\end{command}
\begin{command}
FROM python:3.11-slim-bookworm

ENV PYTHONDONTWRITEBYTECODE=1 \
    PYTHONUNBUFFERED=1 \
    PIP_DISABLE_PIP_VERSION_CHECK=1

RUN apt-get update \
    && DEBIAN_FRONTEND=noninteractive apt-get install -y --no-install-recommends \
        texlive \
        texlive-latex-extra \
        texlive-luatex \
        texlive-xetex \
        texlive-lang-japanese \
        texlive-pictures \
        texlive-science \
        texlive-fonts-extra \
        fonts-noto-cjk \
        fonts-ipaexfont-mincho \
        fonts-ipaexfont-gothic \
    && mktexlsr \
    && updmap-sys \
    && rm -rf /var/lib/apt/lists/*

WORKDIR /app

COPY requirements.txt ./
RUN python -m pip install --no-cache-dir -r requirements.txt

COPY app.py ./

RUN useradd --system --uid 10001 --create-home texnote \
    && chown -R texnote:texnote /app

USER texnote

EXPOSE 8000

HEALTHCHECK --interval=30s --timeout=5s --start-period=30s --retries=3 \
  CMD python -c "import urllib.request; urllib.request.urlopen('http://127.0.0.1:8000/health', timeout=4)" || exit 1

CMD ["uvicorn", "app:app", "--host", "0.0.0.0", "--port", "8000"]
\end{command}

\texttt{app.py}はTeXエンジンを\texttt{/usr/bin}の絶対パスで呼び出す。
Debianパッケージをコンテナへ導入しているため、既存コードを変更せず利用できる。

\subsection{ビルドコンテキストから秘密情報を除外}

\texttt{/home/mat/texnote-server-docker/.dockerignore}を作成する。
\begin{command}
cd /home/mat/texnote-server-docker
nano .dockerignore
\end{command}
\begin{command}
.git
.DS_Store
.venv
__pycache__
*.pyc
.api-token
server.env
docs
out
\end{command}

\texttt{server.env}を除外することが重要である。
これにより、APIトークンがイメージのlayerやbuild contextへ入らない。

\subsection{Compose設定の作成}

\texttt{/home/mat/texnote-server-docker/compose.yaml}を作成する。
\begin{command}
cd /home/mat/texnote-server-docker
nano compose.yaml
\end{command}
\begin{command}
services:
  texnote-server:
    build:
      context: .
      dockerfile: Dockerfile
    image: texnote-server:bookworm
    restart: unless-stopped
    env_file:
      - ./server.env
    ports:
      - "192.168.3.11:8000:8000"
    read_only: true
    tmpfs:
      - /tmp:rw,noexec,nosuid,size=2g,mode=1777
    security_opt:
      - no-new-privileges:true
    cap_drop:
      - ALL
    pids_limit: 128
\end{command}

このAPIは版組ごとに\texttt{/tmp}下へ一時ディレクトリを作り、処理後に削除する。
したがって永続volumeは不要である。
\texttt{read\_only}と\texttt{tmpfs}により、イメージ本体は変更せず一時領域だけを書き込み可能にする。

設定と秘密情報の権限を確認する。
\begin{command}
cd /home/mat/texnote-server-docker
chmod 600 server.env
ls -la
sudo docker compose config --quiet
\end{command}
\texttt{docker compose config}を引数なしで実行すると、展開後の環境変数が表示される場合がある。
トークンを画面やログへ出したくない場合は\texttt{--quiet}だけを使う。

\subsection{イメージのビルドとTeX環境の検査}

systemd版を動かしたままイメージをビルドできる。
初回はTeX Live一式をダウンロードするため時間がかかる。
\begin{command}
cd /home/mat/texnote-server-docker
df -h /
sudo docker compose build
sudo docker image ls texnote-server:bookworm
sudo docker system df
\end{command}

APIサーバーを公開する前に、完成したイメージ内を検査する。
\begin{command}
sudo docker compose run --rm --no-deps texnote-server \
  sh -lc 'uname -m && python --version'

sudo docker compose run --rm --no-deps texnote-server \
  sh -lc 'command -v lualatex; command -v xelatex; command -v pdflatex; command -v uplatex; command -v platex; command -v dvipdfmx; command -v extractbb'

sudo docker compose run --rm --no-deps texnote-server \
  sh -lc 'kpsewhich ipaexm.ttf; kpsewhich ipaexg.ttf; kpsewhich ipxm-r-u61.tfm; kpsewhich ipxm-r-u61.pfb; kpsewhich ipaex-type1.map'
\end{command}
すべての実行ファイルとフォントでパスが表示されることを確認する。

\subsection{systemd版からDocker版への切替}

まずsystemd版の現在状態を記録する。
\begin{command}
systemctl is-enabled texnote-server
systemctl is-active texnote-server
\end{command}
Docker版とポートが競合しないよう、systemd版を停止して自動起動を無効にする。
サービスファイルやPython仮想環境は、切替確認が終わるまで削除しない。
\begin{command}
sudo systemctl disable --now texnote-server
systemctl is-active texnote-server
\end{command}

Docker版をバックグラウンド起動する。
\begin{command}
cd /home/mat/texnote-server-docker
sudo docker compose up -d
sudo docker compose ps
sudo docker compose logs --tail=100 texnote-server
\end{command}

Pi上で確認する。
\begin{command}
curl --fail-with-body http://127.0.0.1:8000/health
sudo docker inspect \
  --format '{{.State.Status}} {{if .State.Health}}{{.State.Health.Status}}{{end}}' \
  texnote-server-docker-texnote-server-1
\end{command}
Composeのproject名やディレクトリ名を変えた場合、コンテナ名も変わる。
正確な名前は\texttt{sudo docker compose ps}で確認する。

MacからLAN経由で確認する。
\begin{command}
curl --fail-with-body http://192.168.3.11:8000/health

TEXNOTE_TOKEN="$(ssh mat@192.168.3.11 \
  "sed -n 's/^TEXNOTE_API_TOKEN=//p' /home/mat/texnote-server-docker/server.env")"
curl --fail-with-body \
  http://192.168.3.11:8000/v1/auth-check \
  -H "Authorization: Bearer $TEXNOTE_TOKEN"
unset TEXNOTE_TOKEN
\end{command}
成功時は\texttt{\{"status":"ok"\}}が返る。
続いてiPad版・iPhone版から実際に5種類のエンジンを順に版組する。

\subsection{再起動後の自動復旧試験}

Composeの\texttt{restart: unless-stopped}とDockerサービスの自動起動を確認する。
\begin{command}
sudo reboot
\end{command}
Piが戻った後、Macから確認する。
\begin{command}
ping -c 3 192.168.3.11
curl --connect-timeout 10 --fail-with-body \
  http://192.168.3.11:8000/health
ssh mat@192.168.3.11 \
  'sudo systemctl is-active docker; cd /home/mat/texnote-server-docker && sudo docker compose ps'
\end{command}

\subsection{更新、ログ、停止}

\subsubsection{Docker Compose版の更新}

\texttt{app.py}または\texttt{requirements.txt}を更新した場合は、
まずPi上の現行ファイルをバックアップする。
\begin{command}
ssh mat@192.168.3.11
cd /home/mat/texnote-server-docker
cp -p app.py app.py.backup
cp -p requirements.txt requirements.txt.backup
exit
\end{command}
Macから更新ファイルを転送する。
\begin{command}
scp \
  /Users/mat/Documents/Git/TeXNote/Server/P2/progs/app.py \
  /Users/mat/Documents/Git/TeXNote/Server/P2/progs/requirements.txt \
  mat@192.168.3.11:/home/mat/texnote-server-docker/

\end{command}
Pi上でPython構文とCompose設定を検査してから、イメージを再ビルドする。
単なる\texttt{docker compose restart}ではイメージ内の\texttt{app.py}が更新されないため、
必ず\texttt{--build}を付ける。
\begin{command}
ssh mat@192.168.3.11
cd /home/mat/texnote-server-docker
python3 -m py_compile app.py
sudo docker compose config --quiet
sudo docker compose up -d --build
sudo docker compose ps
sudo docker compose logs --tail=100 texnote-server
curl --fail-with-body http://192.168.3.11:8000/health
\end{command}
異常があればバックアップを戻し、旧イメージを再ビルドする。
\begin{command}
cd /home/mat/texnote-server-docker
cp -p app.py.backup app.py
cp -p requirements.txt.backup requirements.txt
sudo docker compose up -d --build
sudo docker compose ps
\end{command}

添付リソースAPIを\texttt{fileName}から\texttt{relativePath}へ変更した版では、
Docker版の\texttt{/health}確認後に対応するiPad・iPhone版アプリを更新する。
更新後はサブフォルダ内のリソースを含むCardで版組、保存、再読込を確認する。

日常の確認には次を使う。
\begin{command}
cd /home/mat/texnote-server-docker
sudo docker compose ps
sudo docker compose logs --tail=100 texnote-server
sudo docker compose logs -f texnote-server
sudo docker stats --no-stream
sudo docker system df
\end{command}
ログ追跡はCtrl-Cで終了してよい。コンテナ自体は停止しない。

明示的に停止・再開する場合：
\begin{command}
cd /home/mat/texnote-server-docker
sudo docker compose stop
sudo docker compose start
\end{command}
コンテナとCompose用networkを削除するが、ビルド済みイメージを残す場合：
\begin{command}
cd /home/mat/texnote-server-docker
sudo docker compose down
\end{command}
未使用イメージだけを整理する場合は、削除対象を先に確認する。
\begin{command}
sudo docker image ls
sudo docker image prune
\end{command}
\texttt{docker system prune -a}は広範囲のイメージやcacheを削除するため、
本手順では使用しない。

\subsection{問題発生時にsystemd版へ戻す}

Docker版で問題が発生した場合は、次の順で元の構成へ戻す。
\begin{command}
cd /home/mat/texnote-server-docker
sudo docker compose down
sudo systemctl enable --now texnote-server
systemctl is-enabled texnote-server
systemctl is-active texnote-server
curl --fail-with-body http://192.168.3.11:8000/health
\end{command}
Docker版とsystemd版の両方を同時に起動しない。

\subsection{Docker版完了チェックリスト}
\begin{itemize}
\item[$\square$] Docker EngineとCompose pluginが導入された。
\item[$\square$] Dockerサービスがenabledかつactiveになった。
\item[$\square$] \texttt{server.env}がbuild contextから除外され、権限600になった。
\item[$\square$] イメージ内で5種類のTeXエンジンを発見できた。
\item[$\square$] IPAexとType 1のTFM/PFB/mapを発見できた。
\item[$\square$] systemd版を停止してからDocker版を起動した。
\item[$\square$] \texttt{/health}と\texttt{/v1/auth-check}が成功した。
\item[$\square$] iPad版・iPhone版から実際に版組できた。
\item[$\square$] Pi再起動後にコンテナが自動復旧した。
\item[$\square$] TCP 8000をインターネットへ公開していない。
\item[$\square$] systemd版への復旧手順を残している。
\end{itemize}

\subsection{Docker公式資料}
\begin{itemize}
\item DebianへのDocker Engine導入：
\href{https://docs.docker.com/engine/install/debian/}{Docker Docs: Install Docker Engine on Debian}
\item Linux版Compose plugin：
\href{https://docs.docker.com/compose/install/linux/}{Docker Docs: Install the Compose plugin}
\item Composeの\texttt{env\_file}：
\href{https://docs.docker.com/compose/how-tos/environment-variables/set-environment-variables/}{Docker Docs: Set environment variables}
\item Composeサービス設定：
\href{https://docs.docker.com/reference/compose-file/services/}{Docker Docs: Define services}
\end{itemize}

\section*{補足}
OSのメジャー更新、ユーザー名、IPアドレス、TeXNoteの保存場所を変えた場合は、
本文中のパスと設定値を一括して読み替えること。

\section*{ソース・プログラム}

\lstinputlisting[caption = requirements.txt ,label = text]{progs/requirements.txt}

\lstinputlisting[caption = app.py ,label = app]{progs/app.py}

\end{document}
